\section{怎样辨析文言虚词的一般用法}

大家知道，实词是能够单独用来回答问题、有比较实在的意义的词；虚词则是不能单独用来回答问题、也没有实在的意义，而是具有帮助造句作用的词。一般以名词、动词、形容词、数词、量词和代词为实词，以副词、介词、连词、助词和叹词为虚词。这一点古今汉语基本相同。不同之处，如古汉语中许多代词不能单独回答问题（如“其”“之”等），所以可归入文言虚词。文言虚词数量多，且在字面、用法和作用上与 modern汉语虚词不同，掌握不易，但它在古代汉语语法中作用重要，不能忽视。以下方法可供初学者参考。

1. 先掌握使用频率高、古今差别大的二十余个文言虚词，结合成段成篇的阅读，分类了解各虚词的用法及讲法，再放入文章中体会。反复几次，即可理解重要虚词的语法作用及相应现代汉语含义，为准确解释词义、翻译语句打基础。重要虚词包括：之、其、者、所、于、以、为、与、则、而、虽、然、何、孰、乃、也、矣、焉、耳、乎、哉、夫、盖等。应先掌握其主要用法，以少御繁，逐步扫除阅读障碍。

2. 根据文言虚词在句中的位置及前后词语，辨析它的一般用法。因为有不少文言虚词由于在句子中的位置和它前后词语的不同，用法就不一样。这是掌握一部分文言虚词用法的一个好办法。例如：

\noindent
[\ 之\ ]
\begin{enumerate}[label=(\arabic*)]
	\item 用在动词后，作宾语，代词，相当于“他（们）”“它（们）”。
	\begin{enumerate}[label=\textcircled{\arabic*}]
		\item 陈胜佐之，并杀两尉。（他，指吴广）
		\item 屠暴起，以刀劈狼首，数刀毙之。（它，指狼）
		\item 以吾酌油知之。（它，指手熟能善射的道理）
	\end{enumerate}
	\item 用在定语和中心词间，结构助词，相当于“的”。
	\begin{enumerate}[label=\textcircled{\arabic*}]
		\item 永之人争奔走焉。
		\item 食者甚众，是天下之大残也。
	\end{enumerate}
	\item 用在主语和谓语间，结构助词，取消句子独立性，不译。
	\begin{enumerate}[label=\textcircled{\arabic*}]
		\item 比吾乡邻之死则已后矣。
		\item 孤之有孔明，犹鱼之有水也。
	\end{enumerate}
	\item 用在不及物动词或形容词后，助词，无义，调整音节。
	\begin{enumerate}[label=\textcircled{\arabic*}]
		\item 公将鼓之。刿曰：“未可。”（《曹刿论战》）
		\item 陈涉少时，尝与人佣耕，辍耕之垄上，怅恨久之。（《陈涉世家》）
	\end{enumerate}
	\item 处在谓语位置，动词，相当于“到······去。”
	\begin{enumerate}[label=\textcircled{\arabic*}]
		\item 吾欲之南海。
		\item 辍耕之垄上。
		\item 项伯乃夜驰之沛公军。
	\end{enumerate}
	\item 处在中心词的定语位置，近指代词，相当于“这、这些”。
	\begin{enumerate}[label=\textcircled{\arabic*}]
		\item 均之二策，宁许已负秦曲。（这）（《廉颇蔺相如列传》）
		\item 齐明、周最、陈轸、召滑、楼缓、翟景、苏厉、乐毅之徒通其意。（这些人）（《过秦论》）
	\end{enumerate}
	\item 处在指同类事物的两名词间，近指代词，相当于“这样、这种”。
	\begin{enumerate}[label=\textcircled{\arabic*}]
		\item 公输盘为楚造云梯之械。（这种）（《公输》）
		\item 以君之力，曾不能损魁父之丘。（这样的）（《愚公移山》）
	\end{enumerate}
\end{enumerate}

\noindent
[\ 夫\ ]
\begin{enumerate}[label=(\arabic*)]
	\item 句首发语助词，引起议论，含“大家知道”意味，可不译。
	\begin{enumerate}[label=\textcircled{\arabic*}]
		\item 夫赵强而燕弱。而君幸于赵王，故燕王欲结于君。（《廉颇蔺相如列传》）
		\item 夫精念存想。或泄于目，或泄于口，或泄于耳。（《订鬼》）
	\end{enumerate}
	\item 句中，除实词名词外，为远指代词，相当于“那、那个、那些”，语意较轻。
	\begin{enumerate}[label=\textcircled{\arabic*}]
		\item 予观夫巴陵胜状。在洞庭一湖。（那）（《岳阳楼记》）
		\item 故为之说，以俟夫观人风者得焉。（那些）（《捕蛇者说》）
	\end{enumerate}
	\item 句尾助词，表感叹，相当于“呀、啊”。
	\begin{enumerate}[label=\textcircled{\arabic*}]
		\item 一人飞升，仙及鸡犬，信夫！（《促织》）
		\item 悲夫！有如此之势，而为秦人积威之所劫······（《六国论》）
	\end{enumerate}
\end{enumerate}

\noindent
[\ 焉\ ]
\begin{enumerate}[label=(\arabic*)]
	\item 动词前，疑问代词，表反问，相当于“哪里”。
	\begin{enumerate}[label=\textcircled{\arabic*}]
		\item 非再至，焉知其奇若此？（《游黄山记》）
		\item 且焉置土石？（《愚公移山》）
	\end{enumerate}
	\item 副词或副词性短语后，表句中停顿，无义；或作形容词词尾，相当于“然”（的、地）。
	\begin{enumerate}[label=\textcircled{\arabic*}]
		\item 于是焉河伯欣然自喜。（表示句中的停顿）（《望洋兴叹》）
		\item 盘盘焉，囷囷焉，蜂房水涡。（······地）（《阿房宫赋》）
		\item 徐徐焉实狼其中。（······地）（《中山狼传》）
	\end{enumerate}
	\item 句末，等于介词“于”加代词“之”，有时仅作代词“之”，或表指点注意的语气（呢）。
	\begin{enumerate}[label=\textcircled{\arabic*}]
		\item 吾舅死于虎，吾夫又死焉（“于之”，于虎）。（《礼记·檀弓下》）
		\item 积水成渊，蛟龙生焉（“于之”，在这里）。（《劝学》）
		\item 率妻子邑人，来此绝境，不复出焉（“于之”，从这里）。（《桃花源记》）
		\item 谨食之，时而献焉（“之”，代蛇）。（《捕蛇者说》）
		\item 句读之不知，惑之不解，或师焉，或不焉（呢）。（《师说》）
	\end{enumerate}
\end{enumerate}

\noindent
[\ 其\ ]
\begin{enumerate}[label=(\arabic*)]
	\item 作定语，第三人称代词，相当于“他（她、它的）”、“他（她、它）们的。”
	\begin{enumerate}[label=\textcircled{\arabic*}]
		\item 其剑自舟中坠于水（他的）。（《刻舟求剑》）
		\item 吾视其辙乱，望其旗靡，故逐之（他们的，代“齐国军队的”。）（《曹刿论战》）
		\item 盖予所至，比好游者尚不能十一，然视其左右，来而记之者已少。（它的，代“所至山洞的”。）（《游裹禅山记》）
		\item 徐而察之，则山下皆石穴罅，不知其浅深（它们的，代“石穴罅的”。）（《石钟山记》）
	\end{enumerate}
	\newpage
	\item 作主谓词组或分句主语（不作独立句主语），相当于“他”、“他们”，有时指代说话人或听话人。
	\begin{enumerate}[label=\textcircled{\arabic*}]
		\item 操蛇之神闻之，惧其不已也，告之于帝。（他，代愚公。“其不已”是主谓词组，作“惧”的宾语。）（《愚公移山》）
		\item 其不皆是，因不用之，是不世之巧无由出也。（他，代马钧，作第一个分句的主语。）（《马钧传》）
		\item 倘能屈威，诚副其所望。（我，代说话人周瑜自已。）（《赤壁之战》）
		\item 老臣以媪为长安君计短也，故以为其爱不若燕后。（您，指听话人赵太后。）（《触龙说赵太后》）
	\end{enumerate}
	\item 在句子中作定语用，具有明显的指代意义，是指示代词，相当于现代的“那”、“那个”、“那种”，有时指代许多中的一个，相当于现代汉语的“其中的”。
	\begin{enumerate}[label=\textcircled{\arabic*}]
		\item 既其出，则或咎其欲出者。（那个）（《游褒禅山记》）
		\item 有蒋氏者，专其利三世矣。（那种）（《捕蛇者说》）
		\item 寺僧使小童持斧，于乱石间，择其一二扣之。（其中的）（《石钟山记》）
		\item 蜀之鄙有二僧，其一贫，其一富。（其中的）（《为学一首示子侄》）
	\end{enumerate}
	\newpage
	\item 用在反问句、感叹句或祈使句的句首，有时放在名词或代词之后，往往是作语气副词。有的跟“乎”呼应，表示反问语气，相当于“岂”，或者只起强化反问语气的作用；有的表示测度的语气，相当于“大概”、“恐怕”；有的表示劝勉、希冀的语气，含有“还是”或“该、可”一类的意思。
	\begin{enumerate}[label=\textcircled{\arabic*}]
		\item 一之谓甚，其可再乎？（岂，难道）（《左传·宫之奇谏假道》）
		\item 其如土石何？（强化反问语气，现代汉语中没有和它相当的词语。）（《愚公移山》）
		\item 天下事已可知，吾其左衽矣！（恐怕，大概）（《肥水之战》）
		\item 与尔三矢，尔其无忘乃父之志！（可、该）（《伶官传序》）
		\item 诸君其筹之！（还是，可）（《方腊起义》）
	\end{enumerate}
\end{enumerate}

我们如果根据文言虚词在句中的位置和它前后的词语来辨析文言虚词的用法，就能解决对一批文言虚词的掌握问题，例如：“诸”这个合音词，用在陈述句子中，相当于“之于”；用在疑问句尾，相当于“之乎”。“乎”字用在陈述句子中，相当于介词“于”；用在疑问句尾，相当于“呢、吗”。其他如“者、也、矣、耳、耶、哉、然、盖······”都可以运用这种办法进行辨析。

3.在文言虚词中，有的可以采用配对比较法来进行辨析。例如“者”和“所”，“于”和“乎”，“则”和“而”，“矣”和“也”，“邪（耶）”和“哉”，等等。这些成对的文言虚词有相同的地方，也有不同的地方，只要加以比较，就能牢牢地掌握住它们的用法。试以“者”和“所”为例说明一下。

“者”和“所”都是结构助词，它们经常跟动词或形容词结合，组成一个具有名词性词组性质的结构，叫做“者字结构”、“所字结构”。在这个意义上，“者”和“所”都有点相当于现代的“的”字。例如：
\begin{enumerate}[label=\textcircled{\arabic*}]
	\item 存者且偷生，死者长已矣。（《石壕吏》）
	\item 京中有善口技者。（《口技》）
	\item 婉贞挥刀奋斫，所当无不披靡。（《冯婉贞》）
	\item 莫如以吾所长攻敌所短。（《冯婉贞》）
\end{enumerate}

“存者”、“死者”、“善口技者”中的“者”，都可译作“的”或“的（人）”；“所当”（面对着，指对战的敌人）、所长（擅长的，“长”活用作动词）、“所短”（欠缺的，“短”活用作动词）中的“所”，都可译作“的”或“的人”。因此，有的语言学家称它们是特殊代词。

不同的是：“所”字处在动词或形容词之前，“者”字处在动词或形容词之后。有时候，“所”和“者”同时出现在动词的前和后，如“所击杀者无虑百十人”（《冯婉贞》）。我们把“所击杀者”只译成“被打死的敌人”就行了。

更为不同的是：“者”字可以放在句中小停顿之前，表示下面将要有所解释，这个助词无义，不译，“所”字则不具有这个功能；而“所”字常常与“以”字连用，表示追究一个“为什么”，或者说明“为了什么”，这是“者”字所不具备的功能。例如：
\begin{enumerate}[label=\textcircled{\arabic*}]
	\item 北山愚公者，年且九十，面山而居。（《愚公移山》）
	\item 开火者，军中发枪之号也。（《冯婉贞》）
	\item 师者，所以传道受业解惑也。（《师说》）
	\item 余叩所以。（《狱中杂记》）
	\item 强秦之所以不敢加兵于赵者，徒以吾两人在也。（《廉颇蔺相如列传》）
\end{enumerate}

句子①②的“者”表示停顿，下面将有所解释，无义，不译。句③的“者”表示停顿，下面将有所解释，无义，不译；“所以”说明“是为了什么”，可以译为“用来······的”。句④的“所以”表示追问“为什么”，可以译为“原因、缘故”。句⑤的“所以”说明“为什么”，可以译为“······的原因”；“者”则表示下面要解释原因，略作停顿。

此外，“者”字用在句末，往往作语气助词用，相当于“似的”，如“言之，貌若甚戚者。”（《捕蛇者说》）而“所”字则常常跟“为”呼应，表示被动语气，如“适为虞人所逐。”（《中山狼传》等等，不再细说。

我们再以“而”和“则”配对比较一下。它们相同的方式，是在连接前后两项，表示承接关系时，大致都和现代的“就、便”相当；在连接意思相反相对的两项，表示转折关系时，又都和现代的“却、但是”相当。例如：
\begin{enumerate}[label=\textcircled{\arabic*}]
	\item 一鼓作气，再而衰，三而竭。（就）（《曹刿论战》）
	\item 天龙闻而下之。（就）（《叶公好龙》）
	\item 虽才高于世，而无骄尚之情。（可是）（《张衡传》）
	\item 未有封侯之赏，而听细说，欲诛有功之人，此亡秦之续耳。（却）（《鸿门宴》）
	\item 木受绳则直。（就）（《劝学》）
	\item 每闻琴瑟之声，则应节而舞。（便）（《促织》）
	\item 北虽貌敬，实则愤怒。（却）（《指南录后序》）
	\item 至则无所用，放之山下。（可是）（《黔之驴》）
\end{enumerate}

不同的是：“而”可以连接不分主次、先后的两项，表示并列关系，相当于现代的“又”，如“（狼）性贪而狠，党豺为虐。”（《中山狼传》）也可以连接前后递进的两项，相当于现代的“而且”，如“君子博学而日参省乎己。”（《劝学》）“则”却只能连接前后紧承连贯的两项，不能连接并列递进的两项。

更为不同的是：“而”如果连接有因果关系的两件事，相当于现代的“因而，于是”，如“表恶其能而不能用也。”（《赤壁之战》）“则”却连接有条件关系的两件事，表示在一定条件下所产生的结果，相当于现代的“那么，就”，如“向吾不为斯役，则久已病矣。”（《捕蛇者说》）甚至“而”可以连接状语和谓语，表示偏正关系，有时可译为“地”，如“吾尝终日而思矣，不如须臾之所学也。”（《劝学》）；有时连接主语和谓语，表示假设一种条件，含有“如果，假若”的意思，如“诸君而有意，瞻予马首可也！”（《冯婉贞》）“则”字统没有这样的用法。但是，“则”字可以用在无对待关系的单句的主语之后，和“即”、“乃”相似，起判断作用，相当于现代的“是，就是”，如“此则岳阳楼之大观也。”（《岳阳楼记》）“而”字却没有这种用法。

其它如“之”与“其”、“也”和“矣”、“夫”和“盖”、“哉”和“乎”等等，我们都可以配对比较，进行辨析，必能使概念明确，印象深刻，从而切实地掌握文言虚词的一般用法。
\newpage